\documentclass[11pt]{article}
\usepackage[utf8]{inputenc}
\usepackage[margin=1in, includefoot]{geometry}
\usepackage{tcolorbox, parskip, amssymb, amsmath, booktabs, listings, xcolor, hyperref, float}

% --- Professional Link Setup ---
\definecolor{darkblue}{rgb}{0, 0, 0.5}
\hypersetup{colorlinks=true, linkcolor=darkblue, citecolor=darkblue, urlcolor=blue}

% --- SQL Code Styling ---
\lstset{
    backgroundcolor=\color{gray!10},
    basicstyle=\ttfamily\small,
    breaklines=true,
    frame=single,
    keywordstyle=\color{blue},
    commentstyle=\color{green!50!black},
    language=SQL
}

\begin{document}

\begin{tcolorbox}[colback=darkblue!5, colframe=darkblue!75!black, title=\centering\Huge{README: myTunes Database Assembly}]
    \LARGE
    \textbf{Course:} COMP 3005AB - Assignment 3 \\
    \textbf{Student:} Rahie Datta (101263083) \\
    \textbf{Instructor:} Dr. Louis D. Nel \hfill \textbf{Date:} \today
\end{tcolorbox}

\section{Project Overview}
This project involved the consolidation of diverse SQL datasets into a unified, normalized \texttt{myTunes} database. Following the \textit{``Practice, Practice, Practice''} philosophy, the implementation ensures 100\% referential integrity and optimized query performance.

\section{Problem 1: Telephone Scenario (P1.01--P1.10)}
The first phase involved querying a telephone system database. Below are the Relational Algebra (RA) specifications for the core logic implemented in the \texttt{P1\_XX.sql} scripts \cite{nel_er_model}.

\subsection{Relational Algebra Specifications}
Let $S$=Subscribers, $L$=Lines, $C$=Calls, $V$=Services, $SS$=Service\_Subscribers, and $T$=Trunks.

\begin{enumerate}
    \item \textbf{P1.01 (Directory):} $\Pi_{name, address, area, office, number} (S \bowtie_{S.portid = L.portid} L)$ \cite{nel_rel_alg}
    \item \textbf{P1.02 (Office Utilization):} $_{area, office} \mathcal{G}_{COUNT(*) \rightarrow num\_sub} (L)$ \cite{nel_rel_alg}
    \item \textbf{P1.03 (Intra-switch Calls):} $\Pi_{name} (S \bowtie_{S.portid = C.orig} (C \bowtie L1 \bowtie L2))$ \cite{nel_rel_alg}
    \item \textbf{P1.04 (No Origination):} $\Pi_{name}(S) - \Pi_{name}(S \bowtie_{S.portid = C.orig} C)$ \cite{nel_rel_alg}
    \item \textbf{P1.05 (Total Inactivity):} $\Pi_{name}(S) - \Pi_{name}(S \bowtie_{S.orig} C) - \Pi_{name}(S \bowtie_{S.term} C)$ \cite{nel_rel_alg}
    \item \textbf{P1.06 (Power Users/Division):} $\Pi_{name, service}(S \bowtie SS) \div \Pi_{scode}(V)$ \cite{nel_rel_alg}
    \item \textbf{P1.07 (Trunk Load):} $_{foreign\_sw} \mathcal{G}_{COUNT(call\_id)} (T \bowtie_{T.portid=C.orig \lor T.portid=C.term} C)$ \cite{nel_rel_alg}
    \item \textbf{P1.08 (Service Rank):} $\tau_{sub\_count \downarrow} (_{service} \mathcal{G}_{COUNT(portid)} (V \bowtie_{left} SS))$ \cite{nel_rel_alg}
    \item \textbf{P1.09 (Area 613 Routing):} $\Pi_{portid, foreign\_sw} (T \bowtie \sigma_{area='613'}(TR))$ \cite{nel_rel_alg}
    \item \textbf{P1.10 (Facility State):} $\Pi_{portid} (\sigma_{state='BUSY'}(CH))$ \cite{nel_rel_alg}
\end{enumerate}


\section{Problem 2: Database Assembly \& ETL}
The assembly phase managed the ingestion of 6 TA-contributed datasets. 

\subsection{Data Integrity \& Orphan Pruning}
A total of \textbf{204 orphans} were initially detected. After merging the data and identifying malformed records from specific contributors, a final pruning step was executed:
\begin{lstlisting}
DELETE FROM master_songs WHERE cd_id NOT IN (SELECT cd_id FROM master_cds);
\end{lstlisting}
This resulted in a final count of \textbf{5,532 valid song records} and \textbf{0 orphans}.

\subsection{Final Database Metrics}
\begin{table}[H]
    \centering
    \begin{tabular}{lr}
        \toprule
        \textbf{Metric} & \textbf{Value} \\
        \midrule
        Total Songs Loaded & 5,532 \\
        Total CDs Recorded & 441 \\
        Average Songs/CD & 12.54 \\
        Integrity Status & 100\% Verified \\
        \bottomrule
    \end{tabular}
\end{table}

\section{Maintenance \& Optimization}
The final \texttt{myTunes.db} was optimized using:
\begin{itemize}
    \item \textbf{VACUUM:} To reclaim storage and reduce fragmentation.
    \item \textbf{ANALYZE:} To update the query planner's internal statistics for better join performance.
\end{itemize}

\section{References}
\begin{thebibliography}{9}
    \raggedright
    \bibitem{sqlite_doc} 
    SQLite Consortium, ``SQLite Documentation,'' \textit{sqlite.org}, 2025. 
    Available: \url{https://www.sqlite.org/docs.html}

    \bibitem{nel_rel_alg} 
    L. D. Nel, ``01 Relational Algebra Part 1 Rev1,'' Course Notes for COMP 3005, Carleton University, 2021.
    Available: \url{https://people.scs.carleton.ca/~ldnel/3005winter2026/notes/05%20Query%20Models/01%20Relational%20Algebra%20Part%201%20Rev1.pdf} \\
    \textit{Used for: Formalizing SQL queries P1.01--P1.10 into Relational Algebra notation.}

    \bibitem{nel_er_model} 
    L. D. Nel, ``04 Telephone Switch ER Model Rev 4,'' Course Notes for COMP 3005, Carleton University, 2021. 
    Available: \url{https://people.scs.carleton.ca/~ldnel/3005winter2026/notes/02B%20Sample%20SQLite%20Databases/04%20Telephone%20Switch%20ER%20Model%20Rev%204.pdf} \\
    \textit{Used for: Identifying entity relationships (Subscribers, Lines, Trunks) and schema constraints.}
\end{thebibliography}

\end{document}